\section{Definition}

In this section we define the parameters and terms
that we use to formally describe a bookmark tree.

A bookmark tree, denoted as $T$, is a rooted tree. Each node 
of the tree, denoted as $N$, is
assigned with a label $\lambda$. 
A node without children is a leaf representing a document,
or it is an internal node representing a category. 
A category can contain both child categories and documents.
(In this paper we use the terms
{\it subcategory} when we stress that
a category is contained in another category. Nevertheless, 
all categories are internal nodes of the tree.)
To identify a node uniquely,
we concatenate the labels from the root to that node.
Figure~\ref{fig-labels} shows an example bookmark tree.
The labeled rectangles are categories and shaded
boxes are bookmarks. For example, following a 
label ``News:Stories'' we can locate the 
{\it Stories} subcategory under the {\it News} category.

\begin{figure}[t]
\epsfysize=180pt
\epsfxsize=320pt
\centerline{\epsffile{fig-labels.eps}}
\caption{A Bookmark Tree}
\label{fig-labels}
\end{figure}

The length of the path from the root to a node is the {\it depth}
of the node, denoted as $\delta$. 
The largest depth of any node is the {\it height}
of $T$, denoted as $h$. The number of children of a node,
both subcategories and documents, is 
called the {\it degree} of the node, denoted as $d$. 
When the degree of a category exceeds the splitting
threshold, the category may be split into subcategories.
When the degree of a subcategory becomes too small,
it may be merged back to its parent node, if the
parent node has adequate space. 

%For instance, the depth and degree of node A in Figure~\ref{fig-labels}
%is two and eight (one subcategory and seven documents),
%respectively; the height of the tree is four.
%A directory has a unique name since each tree path is unique.
%However, since a bookmark can belong to multiple
%categories and be stored in
%more than one directory (discussed in
%Section~\ref{sec-get-labels}), aliases of a bookmark may exist.
%We thus maintain an {\it alias table} in which the aliases of
%bookmarks are stored. When a bookmark is deleted by the user,
%all the bookmark's aliases must be removed 
%from the category tree as well as from the alias table.

To assist the reader, Table~\ref{table-parameters}
summarizes these parameters, together with other parameters that will
be introduced later. 

\sspace
\begin{table}
\centerline{
\begin{tabular}{|c|l|} \hline
{\it Parameter} & {\it Description} \\ \hline \hline
{\it $T$} & The bookmark tree \\
{\it $N$} & Total number of bookmarks \\
{\it $h$} & The height of $T$ \\
{\it $V$} & A vertex of $T$ \\
{\it $V.L$} & The label of vertex $V$ \\ 
{\it $V.\delta$} & The depth of vertex $V$ \\ 
{\it $V.d$} & The degree of vertex $V$ \\ 
{\it $V.v$} & The number of children of vertex $V$ \\ 
{\it $V.B$} & The name of the bookmark stored in the leaf vertex \\
{\it $\alpha$} & The maximum degree allowed in a directory \\             
{\it $\delta_{split}$} & The threshold beyond $\alpha$ to split a directory\\
{\it $\delta_{merge}$} & The threshold below $\alpha$ to merge subdirectories \\
{\it $\kappa$} & The maximum allowable aliases of a bookmark \\             
{\it $c_{min}$} & The minimum required coverage of a bookmark label \\      
{\it $g_{min}$} & The minimum required coverage of a group of bookmark labels\\             
\hline
\hline
\end{tabular}}
\caption{Parameters}
\label{table-parameters}
\end{table}
\nspace
                 

